\documentclass[aspectratio=169,xcolor=dvipsnames]{beamer}
\usetheme{SimpleDarkBlue}

\usepackage{hyperref}
\usepackage{graphicx}
\usepackage{booktabs}
\usepackage{amssymb}
\usepackage{amsmath}
\usepackage{tikz}
\usepackage{xcolor}
\usepackage{caption}
\usepackage{amsthm}

\title{Introduction to Cyclic Codes}

\author{
	Mohammad Hosein Rabiee
	\and
	Amir Hosein Karimzadgan \\
}

\institute
{
	Department of Computer Science and Statistics \\
	Faculty of Mathematics \\
	K. N. Toosi University of Technology
}
\date{}

\begin{document}
	
	\begin{frame}
		\titlepage
	\end{frame}
	
	\begin{frame}{Cyclic Codes}
		
		\begin{itemize}
			\item \textcolor{blue}{Cyclic codes} are an important class of linear error-correcting codes
			with strong algebraic structure and efficient implementation via \textcolor{blue}{shift registers}.
		\end{itemize}
		
		\vspace{0.6em}
		
		\begin{block}{Definition (Cyclic Code)}
			A code $C$ of length $n$ over $\mathbb{F}_q$ is \alert{cyclic} if:
			\begin{enumerate}
				\item $C$ is a \textcolor{blue}{linear code}
				\item For every codeword $(a_0,a_1,\dots,a_{n-1}) \in C$,  
				its cyclic shift $(a_{n-1},a_0,\dots,a_{n-2})$ is also in $C$
			\end{enumerate}
			
		\end{block}
		\begin{itemize}
			\item When studying cyclic codes, we label coordinates as
			$0,1,\ldots,n-1$ and associate each vector
			\[
			(a_0,a_1,\ldots,a_{n-1}) \in V(n,q)
			\]
			with the polynomial
			\[
			a_0 + a_1 x + \cdots + a_{n-1} x^{n-1}.
			\]
		\end{itemize}
		
	\end{frame}
	
	\begin{frame}{Cyclic Codes: Examples}
		
		\begin{alertblock}{Examples}
			\begin{itemize}
				\item The binary code
				\[
				\{000,\;101,\;011,\;110\}
				\]
				is \alert{cyclic}, since every cyclic shift of a codeword is again in the code.
				
				\vspace{0.4em}
				
				\item  The Hamming code $\text{Ham}(3,2)$ is \alert{cyclic}.  
				Each codeword is obtained as a cyclic shift of its predecessor.
				
				\vspace{0.4em}
				
				\item The binary linear code
				\[
				\{0000,\;1001,\;0110,\;1111\}
				\]
				is \alert{not cyclic}.  
				However, by interchanging the third and fourth coordinates, it becomes the cyclic code
				\[
				\{0000,\;1010,\;0101,\;1111\}.
				\]
				
			\end{itemize}
		\end{alertblock}
		
	\end{frame}
	
	\begin{frame}{Polynomials over Finite Fields}
		
		\begin{block}{Polynomials over $\mathbb{F}_q$}
			Let $\mathbb{F}_q$ (or simply $\mathbb{F}$) denote a finite field.
			\begin{itemize}
				\item $\mathbb{F}[x]$ denotes the set of polynomials in $x$ with coefficients in $\mathbb{F}$.
				\item For
				\[
				f(x) = f_0 + f_1 x + \cdots + f_m x^m, \quad f_m \neq 0,
				\]
				the integer $m$ is called the \alert{degree} of $f(x)$, denoted $\deg f(x)$.
				\item $f_m$ is the \alert{leading coefficient}.
				\item A polynomial is called \alert{monic} if its leading coefficient is $1$.
			\end{itemize}
		\end{block}
		
		\begin{itemize}
			\item Polynomials in $\mathbb{F}[x]$ can be added, subtracted, and multiplied.
			\item $\mathbb{F}[x]$ is a \textcolor{blue}{ring} (but not a field), since polynomials of degree $>0$ have no multiplicative inverses.
			\item For $f(x), g(x) \in \mathbb{F}[x]$,
			\[
			\deg(f(x)g(x)) = \deg f(x) + \deg g(x).
			\]
		\end{itemize}
		
	\end{frame}
	
	\begin{frame}{Division Algorithm for Polynomials}
		
		\begin{block}{Division Algorithm}
			For any $a(x), b(x) \in \mathbb{F}[x]$ with $b(x) \neq 0$, there exist
			unique polynomials $q(x)$ and $r(x)$ such that
			\[
			a(x) = q(x)b(x) + r(x),
			\quad \deg r(x) < \deg b(x).
			\]
			This is analogous to the division algorithm in $\mathbb{Z}$.
		\end{block}
		
		\vspace{0.4em}
		
		\begin{alertblock}{Example}
			In $\mathbb{F}[x]$, divide $x^3 + x + 1$ by $x^2 + x + 1$:
			\[
			x^3 + x + 1 = (x + 1)(x^2 + x + 1) + x.
			\]
			Here, $q(x) = x + 1$ and $r(x) = x$.
		\end{alertblock}
		
	\end{frame}
	
	\begin{frame}{Polynomials Modulo $f(x)$}
		
		\begin{block}{Congruence Modulo a Polynomial}
			Let $f(x) \in \mathbb{F}[x]$ be fixed.  
			Two polynomials $g(x), h(x) \in \mathbb{F}[x]$ are said to be
			\alert{congruent modulo $f(x)$}, written
			\[
			g(x) \equiv h(x) \pmod{f(x)},
			\]
			if $g(x) - h(x)$ is divisible by $f(x)$.
		\end{block}
		
		\vspace{0.4em}
		
		\begin{block}{The Ring $\mathbb{F}[x]/(f(x))$}
			\begin{itemize}
				\item By the division algorithm, every $a(x) \in \mathbb{F}[x]$ is congruent
				modulo $f(x)$ to a unique polynomial $r(x)$ with
				\[
				\deg r(x) < \deg f(x).
				\]
				\item $\mathbb{F}[x]/(f(x))$ denotes the set of all such polynomials.
				\item Addition is ordinary polynomial addition.
				\item Multiplication is performed modulo $f(x)$.
			\end{itemize}
		\end{block}
		
	\end{frame}
	
	\begin{frame}{Example: $\mathbb{F}_{2}[x]/(x^2 + x + 1)$}
		\begin{alertblock}{}
			In $\mathbb{F}_{2}[x]/(x^2 + x + 1)$,
			\[
			(x+1)^2 = x^2 + 2x + 1 \equiv x \pmod{x^2 + x + 1}.
			\]
		\end{alertblock}
		\begin{alertblock}{}
			The elements of the ring are
			\[
			\{\,0,\;1,\;x,\;1+x\,\}.
			\]
			Every non-zero element has a multiplicative inverse, hence
			$\mathbb{F}_{2}[x]/(x^2 + x + 1)$ is a \alert{field} of order $4$.
			
			\vspace{0.6em}
			
			\begin{columns}[c]
				\column{0.48\textwidth}
				\centering
				\textbf{Addition ($+$)}
				\[
				\begin{array}{c|cccc}
					+ & 0 & 1 & x & 1+x \\ \hline
					0 & 0 & 1 & x & 1+x \\
					1 & 1 & 0 & 1+x & x \\
					x & x & 1+x & 0 & 1 \\
					1+x & 1+x & x & 1 & 0
				\end{array}
				\]
				
				\column{0.48\textwidth}
				\centering
				\textbf{Multiplication ($\cdot$)}
				\[
				\begin{array}{c|cccc}
					\cdot & 0 & 1 & x & 1+x \\ \hline
					0 & 0 & 0 & 0 & 0 \\
					1 & 0 & 1 & x & 1+x \\
					x & 0 & x & 1+x & 1 \\
					1+x & 0 & 1+x & 1 & x
				\end{array}
				\]
			\end{columns}
			
		\end{alertblock}
		
	\end{frame}
	
	\begin{frame}{Irreducible Polynomials}
		
		\begin{alertblock}{When is $\mathbb{F}[x]/(f(x))$ a Field?}
			It is \alert{not} true that $\mathbb{F}[x]/(f(x))$ is a field for every
			polynomial $f(x)$.
			
			\begin{itemize}
				\item For example, $\mathbb{F}[x]/(x^2 + 1)$ is \alert{not} a field.
				\item The key property that makes $\mathbb{F}[x]/(f(x))$ a field
				is that $f(x)$ is \alert{irreducible}.
			\end{itemize}
		\end{alertblock}
		
		\begin{block}{Definition (Reducible / Irreducible)}
			A polynomial $f(x) \in \mathbb{F}[x]$ is called \alert{reducible} if
			\[
			f(x) = a(x)b(x),
			\]
			where $a(x), b(x) \in \mathbb{F}[x]$ and
			\[
			\deg a(x) < \deg f(x), \quad \deg b(x) < \deg f(x).
			\]
			If $f(x)$ is not reducible, it is called \alert{irreducible}.
		\end{block}
		
	\end{frame}
	
	\begin{frame}{Factorization in $\mathbb{F}[x]$}
		
		\vspace{0.6em}
		
		\begin{block}{Lemma}
			Let $f(x) \in \mathbb{F}[x]$.
			\begin{enumerate}
				\item $f(x)$ has a linear factor $(x-a)$ \alert{if and only if}
				\[
				f(a) = 0.
				\]
				
				\item If $\deg f(x) = 2$ or $3$, then $f(x)$ is \alert{irreducible}
				if and only if
				\[
				f(a) \neq 0 \quad \text{for all } a \in \mathbb{F}.
				\]
				
				\item Over any field,
				\[
				x^n - 1 = (x - 1)\left(x^{n-1} + x^{n-2} + \cdots + x + 1\right),
				\]
				where the second factor may be further reducible.
			\end{enumerate}
		\end{block}
		
	\end{frame}
	
	\begin{frame}{Factorization in $\mathbb{F}[x]$ (Continued)}
		
		\begin{alertblock}{Proof}
			Let $f(x) \in \mathbb{F}[x]$.
			\begin{enumerate}
				\item[(i)] If $f(x) = (x-a)g(x)$, then clearly $f(a) = 0$.
				Conversely, if $f(a) = 0$, then by the division algorithm:
				\[
				f(x) = q(x)(x-a) + r(x), \quad \deg r(x) < 1.
				\]
				Hence $r(x)$ is a constant, and since $0 = f(a) = r(a)$, we have $r(x) = 0$.
				
				\item[(ii)] A polynomial of degree 2 or 3 is reducible \emph{iff} it has at least one linear factor.
				By (i), the result follows immediately.
				
				\item[(iii)] By (i), $x-1$ is a factor of $x^n - 1$.
				Performing long division of $x^n - 1$ by $x-1$ gives the other factor:
				\[
				x^n - 1 = (x-1)(x^{n-1} + x^{n-2} + \cdots + x + 1).
				\]
			\end{enumerate}
		\end{alertblock}
		
	\end{frame}
	
	\begin{frame}{Finite Fields $\mathrm{GF}(p^h),\ h>1$}
		
		\begin{alertblock}{Key Property}
			The property of a polynomial being irreducible in $\mathbb{F}[x]$ corresponds exactly to the property of a number being prime in $\mathbb{Z}$.
		\end{alertblock}
		
		\begin{block}{Theorem}
			The ring $\mathbb{F}[x]/f(x)$ is a field \emph{if and only if} $f(x)$ is irreducible in $\mathbb{F}[x]$.
		\end{block}
		
		\textbf{Remarks:}
		\begin{itemize}
			\item Proof is exactly in same the way of proving $Z_m$ is a field if and only if $m$ is prime.
			\item For any prime $p$ and any positive integer $h$, there exists an irreducible polynomial over $\mathrm{GF}(p)$ of degree $h$.
			\item By Theorem, this guarantees the existence of fields $\mathrm{GF}(p^h)$ for all integers $h \ge 1$.
		\end{itemize}
		
	\end{frame}
	
	\begin{frame}{Back to Cyclic Codes}
		
		\begin{itemize}
			\item We now return to \alert{cyclic codes}, equipped with polynomial tools.
			\item The key idea is to represent vectors by polynomials and cyclic shifts by
			multiplication by $x$.
		\end{itemize}
		
		\begin{block}{Vector–Polynomial Correspondence}
			A vector $(a_0,a_1,\dots,a_{n-1}) \in \mathbb{F}_q^n$ corresponds to
			\[
			a(x) = a_0 + a_1 x + \cdots + a_{n-1} x^{n-1}.
			\]
		\end{block}
		
	\end{frame}
	
	\begin{frame}{The Ring $R_n = \mathbb{F}_q[x]/(x^n - 1)$}
		
		\begin{itemize}
			\item In $R_n$, multiplication is performed modulo $x^n - 1$.
			\item Since $x^n \equiv 1 \pmod{x^n - 1}$, multiplication by $x$
			causes coefficients to wrap around.
		\end{itemize}
		
		\begin{alertblock}{Key Observation}
			Multiplication by $x$ in $R_n$ corresponds exactly to a cyclic shift of a vector.
		\end{alertblock}
		
	\end{frame}
	
	\begin{frame}{Algebraic Characterization of Cyclic Codes}
		
		\begin{theorem}
			A subset $C \subseteq R_n$ is a cyclic code if and only if:
			\begin{enumerate}
				\item $C$ is closed under addition;
				\item for all $a(x) \in C$ and $r(x) \in R_n$, we have $r(x)a(x) \in C$.
			\end{enumerate}
		\end{theorem}
		
		\begin{block}{Interpretation}
			Cyclic codes are exactly the \alert{ideals} of the ring $R_n$.
		\end{block}
		
	\end{frame}
	
	\begin{frame}{Why This Works}
		
		\begin{itemize}
			\item Cyclicity means closure under multiplication by $x$.
			\item Closure under multiplication by $x$ implies closure under any
			polynomial in $R_n$.
			\item Linearity ensures closure under addition.
		\end{itemize}
		
		\begin{alertblock}{Takeaway}
			Cyclic structure becomes simple algebra.
		\end{alertblock}
		
	\end{frame}
	
	\begin{frame}{Generator Polynomial}
		
		\begin{theorem}
			Let $C$ be a nonzero cyclic code in $R_n$. Then:
			\begin{enumerate}
				\item There exists a unique \alert{monic} polynomial $g(x)$ of smallest degree in $C$.
				\item $C = (g(x)) = \{a(x)g(x) \mid a(x) \in R_n\}$.
				\item $g(x)$ divides $x^n - 1$.
			\end{enumerate}
		\end{theorem}
		
	\end{frame}
	
	\begin{frame}{Properties of the Generator Polynomial}
		
		\begin{itemize}
			\item The generator polynomial completely determines the code.
			\item If $\deg g(x) = r$, then
			\[
			\dim(C) = n - r.
			\]
			\item Every cyclic code corresponds to a monic divisor of $x^n - 1$.
		\end{itemize}
		
	\end{frame}
	
	\begin{frame}{Example: Binary Cyclic Codes of Length 3}
		
		Over $\mathbb{F}_2$,
		\[
		x^3 - 1 = (x + 1)(x^2 + x + 1),
		\]
		where both factors are irreducible.
		
		\begin{itemize}
			\item Each monic divisor gives a cyclic code.
			\item This produces all binary cyclic codes of length $3$.
		\end{itemize}
		
	\end{frame}
	
	\begin{frame}{Check Polynomial}
		
		Let $g(x)$ be the generator polynomial of a cyclic code.
		
		\begin{block}{Definition}
			The \alert{check polynomial} $h(x)$ is defined by
			\[
			x^n - 1 = g(x)h(x).
			\]
		\end{block}
		
		\begin{theorem}
			A polynomial $c(x) \in R_n$ is a codeword if and only if
			\[
			c(x)h(x) \equiv 0 \pmod{x^n - 1}.
			\]
		\end{theorem}
		
	\end{frame}
	
	\begin{frame}{Parity-Check Matrix}
		
		\begin{itemize}
			\item The check polynomial encodes all parity conditions.
			\item Writing the coefficients of
			\[
			h(x),\; xh(x),\; x^2h(x), \dots
			\]
			yields the rows of a parity-check matrix.
		\end{itemize}
		
		\begin{alertblock}{Summary}
			$g(x)$ generates the code; $h(x)$ checks membership.
		\end{alertblock}
		
	\end{frame}
	
	\begin{frame}{Final Takeaway}
		
		\begin{itemize}
			\item Cyclic codes turn coding problems into algebraic ones.
			\item Their structure is governed by the factorization of $x^n - 1$.
			\item Generator and check polynomials enable efficient encoding and decoding.
		\end{itemize}
		
		\begin{center}
			\alert{Cyclic codes = structure + efficiency}
		\end{center}
		
	\end{frame}
	
	\begin{frame}
		\Huge{\centerline{\textbf{Thank you!}}}
	\end{frame}
	
\end{document}
